\section{FORMAL RESTATEMENT}
\textbf{Erd\H{o}s \#144; partial reconstruction, 2026-09-23.}
Prove that the positive integers $n$ having divisors $d_1,d_2$ with
$d_1<d_2<2d_1$ have natural density one. The stronger version replaces
$2$ by any fixed $c>1$.

\section{QUICK LITERATURE AND CONTEXT CHECK}
Maier--Tenenbaum proved the assertion, including its stronger version.
Their primary paper concerns the distribution and concentration of
divisors. We give an exact reduction to a set of multiples, a quantitative
sufficient condition, and a proved positive-density family. These do
not establish density one; the missing analytic step is identified
explicitly. No new resolution is claimed.

\section{WORK}
Fix $c>1$. Define
\[
 \mathcal M_c=\{ab:1\le a<b<ca,\ (a,b)=1\}.
\]
Then $n$ has the desired close divisors if and only if it is divisible
by a member of $\mathcal M_c$. For the forward implication, write
$d_1=ga,d_2=gb$ after cancelling their gcd. Since
$\operatorname{lcm}(d_1,d_2)=gab$ divides $n$, so does $ab$.
Conversely, if $ab\mid n$ with the stated inequalities, $a,b$ themselves
are suitable divisors. This proves the exact identity
\[
 \{n:\exists d_1\mid n,\ d_2\mid n,\ d_1<d_2<cd_1\}
 =\bigcup_{m\in\mathcal M_c}m\mathbb N.                 \tag{1}
\]

For any finite $F\subset\mathcal M_c$ the union of its multiples is
periodic of period $L=\operatorname{lcm}(F)$ and has density
\[
 \delta_F=\sum_{\varnothing\ne J\subseteq F}
          \frac{(-1)^{|J|+1}}{\operatorname{lcm}(J)}.      \tag{2}
\]
This is inclusion--exclusion, since the intersection of the sets of
multiples for $J$ is exactly the multiples of its least common multiple.
Thus the lower density in (1) is at least every $\delta_F$.
In particular, it would suffice to find finite sets $F$ with
$\delta_F\to1$. Large overlapping collections cannot be counted by
adding their reciprocal densities; (2) displays the overlaps that
must be controlled.

For $c=2$, take the four coprime pairs $(2,3),(3,5),(4,5),(5,7)$,
so $F=\{6,15,20,35\}$. Among residues modulo $420$, their individual
counts are $70,28,21,12$. Pairwise intersection counts sum to
$14+7+2+7+4+3=37$; triple intersection counts sum to
$7+2+1+1=11$; the fourfold count is $1$. Consequently
\[
 \delta_F=(131-37+11-1)/420=26/105.                       \tag{3}
\]
This proves a concrete positive lower density. For every $c>1$, choose
an integer $a>1/(c-1)$; then $a<a+1<ca$, so all multiples of
$a(a+1)$ supply a positive-density family in the stronger version too.

\subsection*{A logarithmic pigeonhole criterion}
Suppose $m\mid n$ and list the divisors of $m$ as
$1=t_1<t_2<\cdots<t_s=m$. If every consecutive ratio is at least $c$,
then multiplying them gives $m\ge c^{s-1}$. Therefore
\[
 \tau(m)>1+\frac{\log m}{\log c}
 \quad\Longrightarrow\quad
 \text{$n$ has two divisors of ratio strictly between $1$ and $c$}.
                                                               \tag{4}
\]
More quantitatively, since
$\sum_{i=1}^{s-1}\log(t_{i+1}/t_i)=\log m$, at most
$\log m/\log c$ of the consecutive ratios can be at least $c$.
Thus at least $s-1-\lfloor\log m/\log c\rfloor$ are less than $c$
whenever that number is positive. Taking $m=n$ is allowed, but no
argument here proves that (4) holds for almost all $n$ or for a
suitable divisor of almost all $n$.

The distinction between $d_2<2d_1$ and $d_2\le2d_1$ matters. Every
power of two has no pair of divisors with ratio strictly between $1$
and $2$, so there are infinitely many exceptions. They have density
zero, and therefore do not contradict the theorem. The witness pairs
used in (3) satisfy the strict inequality.

\section{VERIFICATION}
The batch script verifies (3) both by direct residue enumeration and
inclusion--exclusion, and checks (1) on finite ranges. Such finite
checks establish the advertised positive-density subfamily only.
A general increasing union of periodic sets need not have a natural
density equal to a naive interchange of limits, so no such interchange
is used to assert the density-one conclusion.

\section{FINAL STATUS}
\textbf{UNRESOLVED within this partial reconstruction; the density-one
theorem is known.}
The exact multiples reduction, density $26/105$ subfamily, and criterion
(4) are proved. The gap is showing that the exceptional set in (1)
is $o(X)$ up to $X$, or an equivalent sufficiently strong divisor-
concentration estimate. A citation to the Maier--Tenenbaum theorem
is not counted as filling that gap in our own argument.

\section{SOURCES}
\url{https://www.erdosproblems.com/144}.
H.~Maier and G.~Tenenbaum, \emph{On the set of divisors of an integer},
Inventiones Mathematicae 76 (1984), 121--128,
\url{https://tenenb.perso.math.cnrs.fr/PPP/Conj-Erdos.pdf}.
Checked 2026-09-23. Completion assesses the limited proof above, not
the established literature status.

\section{COMPLETION ESTIMATE}
\textbf{COMPLETION: 15\%}
